\section{Introduction}\label{sec:introduction}
Traffic forecasting is increasingly framed as a prediction problem on a network, but an uncertainty service also depends on a reporting process. A speed measurement has a physical timestamp, a delivery timestamp and a role in subsequent decisions. If a detector gateway stops reporting, the forecasting model loses recent inputs. The calibration module also loses the realised outcomes needed to assess forecasts already issued. Treating these two losses as interchangeable hides an important distinction: a predictor may continue producing plausible values while the residual evidence supporting its intervals becomes stale. The uncertainty problem therefore concerns information availability as well as traffic dynamics.

The practical motivation is straightforward. An operator receiving a forecast needs to know how much confidence to place in it during degraded sensing. A narrow interval based on a healthy historical period can be misleading after a correlated outage, even when its point forecast changes little. Conversely, widening every interval indiscriminately can make the service uninformative. A useful evaluation must consider coverage together with interval width and the penalty for missed outcomes. It must also retain difficult targets in the scoring population when their reports are unavailable to calibration. Otherwise, removing the very outcomes most affected by failure can artificially improve reported performance.

Delayed feedback presents a second difficulty. A backlog released after restoration may contain many residuals, yet describe conditions several hours earlier. Delivery recency is then a poor proxy for relevance. Permanent loss creates a different information deficit because the missing residuals never enter the live archive. A method that succeeds with delayed packets need not succeed when outcomes disappear. These mechanisms motivate retaining target timestamps, issuing immutable intervals and distinguishing input release from feedback release. The proposed replay follows this ordering explicitly rather than giving the calibrator access to evaluator truth or retrospectively revising uncertainty after the outcome becomes known.

Spatial information offers a possible response, but creates its own risks. Nearby sensors can supply residual evidence when a local detector is silent. Their errors can also be correlated, and their speed regimes need not match the affected location. Context matching and temporal weighting may improve relevance while reducing effective support. An historical archive can stabilise sparse updates, but can preserve obsolete residual behaviour. The methodological question is consequently not whether every available record should be pooled. It is how the usefulness of live, neighbouring and historical evidence changes under a specified reporting process, and how those choices affect interval quality on held-out data.

Traffic predictors provide several ways to represent network dependence. Graph WaveNet and DCRNN use graph-based temporal modelling, whereas STID and STAEformer use identity information and adaptive embeddings \citep{wu2019,li2018,shao2022,liu2023}. The present study does not seek to replace these forecasting architectures. It holds each fitted predictor fixed and compares calibrators on identical forecasts, scales and evaluation masks. This separation makes a paired score difference interpretable as a calibration effect conditional on the backbone. Repeating the comparison across backbones then tests transfer without confusing a stronger point predictor with a better use of outcome feedback.

Conformal and adaptive inference offer relevant foundations, but their assumptions must be respected. ACI adjusts uncertainty from sequential coverage errors, and multistep extensions account for the ordinary wait until a forecast target occurs \citep{gibbs2021,hallberg2024}. Recent work also studies delayed or corrupted feedback directly \citep{elhalabi2026,wang2026}. Accordingly, feedback-aware uncertainty is not claimed as an unexplored general topic. Our narrower focus is a traffic replay with independently impaired input and outcome channels, correlated outages, permanent loss and auditable historical fallback. The proposed weighted signed-residual method is evaluated empirically; source theorems are not transferred to its modified reporting protocol.

The research has three connected objectives. First, it establishes an information boundary that prevents unavailable outcomes from entering calibration while allowing originally valid targets to remain assessable by the evaluator. Second, it evaluates whether asymmetric residual pooling improves the registered interval-score endpoint without falling below a prespecified coverage floor. Third, it investigates whether the result survives changes in network, forecasting backbone, reporting severity and individual components. Each objective requires a different comparison. Fresh simulation episodes assess the locked candidate; matched-rate controls examine an alternative outage geometry; component removals assess mechanisms while retaining the same point and scale caches.

The empirical design combines METR-LA, PEMS-BAY and controlled SUMO simulations. Real-data inference uses chronological partitions and fault-focused hourly origins, while simulation permits full episode evaluation under separately specified physical and reporting events. The study includes core conditions, independent controls, one-factor stress and seven complete component ablations. The breadth is useful only if scope differences remain visible. An hourly network audit is not a complete five-minute deployment trial, and controlled simulation is not proof of field effectiveness. The manuscript therefore reports sampling rules, adaptation choices, uncertainty intervals and adverse findings alongside favourable results, rather than presenting a single aggregated ranking.

The main contribution is an executable causal benchmark coupled with a feedback-aware calibration design and a transparent account of its conditional performance. The final asymmetric method passes twelve of sixteen real-data backbone cells, but this aggregate conceals important differences between networks. Full ablations also reveal components whose removal improves score. Such findings constrain both mechanism claims and practical recommendations. No universal superiority, arbitrary-missingness coverage guarantee or identifiable correction for missing-not-at-random outcomes is asserted. Potential societal benefits motivate the work, while actual impacts on traffic control, safety, travel time or equity remain outside the demonstrated outcomes.

A reproducible comparison also needs recoverable provenance. Saved schedules, forecast identifiers, configuration hashes and originally valid masks allow readers to distinguish a statistical result from an interrupted runner or a reporting artifact. This accountability is also particularly important when evaluations are resumed, because existing verified stages remain evidence rather than being silently replaced.

The remainder of the paper follows the dependency order of the study. Section~\ref{sec:related} positions forecasting, missing-data reconstruction and sequential calibration. Section~\ref{sec:data} describes datasets and partitions. Section~\ref{sec:architecture} explains the information flow. Section~\ref{sec:math} formalises the model, and Section~\ref{sec:algorithm} details its algorithm. Section~\ref{sec:protocol} defines comparisons and inference. Section~\ref{sec:results} presents results, Section~\ref{sec:ablation} examines ablations, and Section~\ref{sec:discussion} discusses interpretation and limitations before Section~\ref{sec:conclusion} concludes.
